\section{Introduction}
\label{sec:introduction}

Policies under hidden dynamics often act before the relevant physical context is known. Two material instances may look identical and still respond differently to contact, load, or motion. Reactive policies map state to action directly; context-conditioned and recurrent policies infer hidden dynamics from experience gathered during the task~\cite{tobin2017domain,peng2018sim2real,hausman2018learning,xie2018learning}. In contact-rich manipulation, that experience is not passive. The push that reveals stiffness, cohesion, or rebound also changes the state the policy must control.

Most rapid-adaptation methods rely on information that appears naturally during the task rollout. That assumption breaks when the useful query differs from the reward-seeking motion, or when the policy must choose a strategy before the dynamics reveal themselves. Explicit probing queries the environment before acting, but the query is itself a physical intervention: the learned latent and the induced state shift jointly determine whether downstream control improves.

\begin{figure}[t]
    \centering
    \includegraphics[width=0.74\textwidth]{figures/fig1_workflow.pdf}
    \caption{\textbf{\PTA{} workflow.} A fixed probe yields a response trace, a paired encoder maps it to $(z,\sigma)$, and a residual PPO policy conditions on $(s_t,z,\sigma)$.}
    \label{fig:pta_workflow}
\end{figure}

We study this tradeoff in cross-material manipulation, where hidden stiffness, cohesion, and flow resistance shape contact-rich control~\cite{huang2022plasticinelab,lin2022diffskill,shi2023robocook}. \PTA{} operationalizes explicit probing (\cref{fig:pta_workflow}) as a controlled context-acquisition instrument: instead of optimizing probe design, it isolates whether an information-gathering interaction remains useful after it has physically perturbed the material.

\paragraph{Recoverability hypothesis.}
Probing is not free: it can displace particles, break contacts, and consume task time. We frame the tradeoff through recoverability. The resulting prediction is directional and falsifiable: probing is advantageous when the elicited response is informative and the induced disturbance relaxes on the task timescale, and costly when the disturbance persists. Elastoplastic media rebound in our benchmark, while sand and snow remain negative under contact loading.

\paragraph{Empirical evaluation.}
We evaluate \PTA{} on a Genesis edge-push benchmark~\cite{genesis2024} spanning granular sand, cohesive snow, elastoplastic material with recoverable rebound, and two parameter-shifted sand variants. A scripted controller has a 55pp transfer spread across the main families, making the splits discriminative; all learned policies train on sand and evaluate zero-shot on the other four splits.

The central result is a material-dependent sign pattern rather than a uniform leaderboard claim. \PTA{} gains 14.7pp on elastoplastic transfer (60.7\% vs.\ 46.0\%) while paying a 13--22pp cost on ID sand, snow, and shifted sand (\cref{tab:main_results,fig:main,fig:effect_delta}). A separate matched-encoder audit strengthens the elastoplastic direction: six matched M7 seeds average 93.5\% transfer versus 39.4\% for six reactive seeds, with a +54.2pp mean paired delta, four wins, and two near ties (\cref{tab:matched_encoder_audit}). The negative splits are evidence of probing's physical cost, not discarded failures. We report effect sizes with per-seed dispersion rather than significance claims.

Ablations clarify the component roles: removing the probe drops elastoplastic transfer below M1, while removing only the encoder output returns near M1-level transfer (\cref{tab:ablation}). A single-seed privileged-parameter baseline (M8) tests ground-truth parameter access under the same sand-only residual-control setup.

\paragraph{Contributions.}
We contribute: (i) a framing of probing as physical intervention; (ii) a five-split Genesis benchmark with sand-only training, per-seed metrics, ablations, and matched-encoder audit; (iii) an empirical signature in which probing is beneficial on recoverable elastoplastic material and costly under persistent particle rearrangement; and (iv) a recoverability hypothesis for additional material families and probe designs.

\paragraph{Scope of claims.}
\PTA{} uses a fixed open-loop probe, simulator-derived aggregate-particle statistics, and an uncalibrated task-conditioning latent. These choices make the claim precise: we identify the phenomenon and recoverability hypothesis in controlled simulation, leaving deployment robustness, learned probes, and sensor calibration as subsequent tests rather than assumptions.
